\documentclass[11pt,a4paper]{article}

\usepackage[utf8]{inputenc}
\usepackage[T1]{fontenc}
\usepackage[french]{babel}
\usepackage{lmodern}
\usepackage{geometry}
\usepackage{amsmath,amssymb,amsfonts}
\usepackage{bm}
\usepackage{graphicx}
\usepackage{booktabs}
\usepackage{hyperref}
\usepackage{xcolor}
\usepackage[most]{tcolorbox}

\geometry{margin=2.5cm}

\title{
\textbf{Couplage gravitationnel de la matière émergente}\\[0.4em]
\large De \(S_{\rm flux}\) au potentiel effectif BuP
}

\author{
Farid Hamdad\\
\small Bottom-Up Quantum Gravity Program
}

\date{2026}

\begin{document}

\maketitle

\begin{abstract}

Paper~8 a construit une source tensorielle matière candidate à partir d'une
excitation locale du réseau d'intrication BuP. Cette source,
\[
S_{\rm flux}
=
T_{00}
-\frac12T_{aa}
+
\frac12T_{\rm grad}
+
\|T_{0a}\|,
\]
combine densité informationnelle, pression effective, contrainte interne et
flux d'intrication.

Le présent papier couple cette source à une limite faible de l'équation
d'Einstein effective BuP. Sur le graphe d'intrication, nous résolvons une
équation de Poisson discrète :
\[
L_{\rm ent}\Phi_{\rm BuP}
=
S_{\rm flux}.
\]
Nous testons ensuite si le potentiel obtenu \(\Phi_{\rm BuP}\) prédit la
réponse de courbure nodale \(|\delta R|\).

Pour \(N=20\), \(\lambda=0.57\), \(k=5\) et \(\sigma=0.15\), nous obtenons
\[
\rho_{\rm Spearman}(S_{\rm flux},|\delta R|)
=
0.741,
\qquad
p=1.84\times10^{-4},
\]
et, après résolution de Poisson,
\[
\rho_{\rm Spearman}(\Phi_{\rm BuP},|\delta R|)
=
0.738,
\qquad
p=2.01\times10^{-4}.
\]
Le potentiel effectif conserve donc presque toute l'information géométrique
portée par la source.

Ce résultat fournit un premier couplage physique entre la matière émergente
et la géométrie effective dans BuP.

\end{abstract}

\tableofcontents

\section{Introduction}

Les papiers précédents ont construit progressivement une chaîne d'émergence.

Paper~7 a montré que les excitations du réseau d'intrication possèdent une
dynamique interne et peuvent former des régimes source-like métastables.

Paper~8 a ensuite construit une source tensorielle matière candidate :
\[
\delta W^{\rm loc}
\longrightarrow
S_{\rm flux}.
\]

La question de Paper~9 est maintenant :

\[
\boxed{
\text{cette source génère-t-elle un potentiel gravitationnel effectif ?}
}
\]

Nous testons donc la chaîne :
\[
\delta W^{\rm loc}
\longrightarrow
S_{\rm flux}
\longrightarrow
\Phi_{\rm BuP}
\longrightarrow
|\delta R|.
\]

\section{Rappel : source tensorielle issue de Paper 8}

Paper~8 a montré que la densité seule \(T_{00}\) est insuffisante pour prédire
la réponse géométrique. La meilleure source testée est :
\[
S_{\rm flux}
=
T_{00}
-\frac12T_{aa}
+
\frac12T_{\rm grad}
+
\|T_{0a}\|.
\]

Les composantes sont :

\[
T_{00}(i)
=
\frac12
\sum_j
(\delta W_{ij}^{\rm loc})^2,
\]

\[
T_{aa}(i)
=
\frac12
\sum_j
(\delta W_{ij}^{\rm loc})^2
\|v_{ij}\|^2,
\]

\[
T_{\rm grad}(i)
=
\sum_j
A_{ij}
\left(
T_{00}(j)-T_{00}(i)
\right)^2,
\]

et :

\[
\vec f(i)
=
\sum_j
(\delta W_{ij}^{\rm loc})^2
\vec v_{ij},
\qquad
\|T_{0a}\|\equiv \|\vec f(i)\|.
\]

Ainsi, la source gravitationnelle effective ne dépend pas seulement d'une
densité scalaire, mais d'une structure tensorielle informationnelle.

\section{Limite faible et équation de Poisson discrète}

L'équation effective introduite dans Paper~5 est :
\[
G_{\mu\nu}[g^{\rm ent}]
=
8\pi G_{\rm eff}(d_s)T_{\mu\nu}^{\rm matter}
+
8\pi G T_{\mu\nu}^{\rm ent}[d_s].
\]

Dans une limite faible, quasi-newtonienne, on cherche un potentiel effectif
\(\Phi_{\rm BuP}\) tel que :
\[
L_{\rm ent}\Phi_{\rm BuP}
=
\mathcal{N}S_{\rm flux}.
\]

Dans les tests numériques, la constante \(\mathcal{N}\) est absorbée dans la
normalisation de la source. Nous résolvons donc :
\[
L_{\rm ent}\Phi_{\rm BuP}
=
S_{\rm flux}.
\]

Le Laplacien \(L_{\rm ent}\) est le Laplacien combinatoire du graphe
d'intrication. Comme il possède un mode zéro, la source est recentrée et une
régularisation faible est utilisée pour résoudre le système.

\section{Observables testées}

Nous comparons trois quantités à la réponse de courbure :

\[
S_{\rm flux},
\qquad
\Phi_{\rm BuP},
\qquad
|\nabla\Phi_{\rm BuP}|.
\]

La réponse géométrique est :
\[
|\delta R_i|
=
|R_i[W^{(0)}+\delta W]-R_i[W^{(0)}]|.
\]

Le critère principal est :
\[
\rho_{\rm Spearman}
\left(
\Phi_{\rm BuP},
|\delta R|
\right).
\]

\section{Protocole numérique}

Le cas principal utilise :
\[
N=20,
\qquad
\lambda=0.57,
\qquad
k=5,
\qquad
A=0.15,
\qquad
\sigma=0.15.
\]

Les poids de la source sont repris de Paper~8 :
\[
\omega=-0.5,
\qquad
\chi=0.5,
\qquad
\psi=1.0.
\]

La source normalisée est donc :
\[
S_{\rm flux}
=
T_{00}
-\frac12T_{aa}
+
\frac12T_{\rm grad}
+
\|T_{0a}\|.
\]

\section{Résultat principal}

Pour \(\sigma=0.15\), nous obtenons :
\[
\rho_{\rm Spearman}(S_{\rm flux},|\delta R|)
=
0.741,
\qquad
p=1.84\times10^{-4}.
\]

Après résolution de :
\[
L_{\rm ent}\Phi_{\rm BuP}
=
S_{\rm flux},
\]
le potentiel obtenu vérifie :
\[
\rho_{\rm Spearman}(\Phi_{\rm BuP},|\delta R|)
=
0.738,
\qquad
p=2.01\times10^{-4}.
\]

Ainsi, le passage par l'équation de Poisson discrète conserve presque toute
l'information géométrique portée par la source.

\begin{table}[htbp]
\centering
\begin{tabular}{l c c}
\toprule
Quantité testée contre \(|\delta R|\) & Spearman & \(p\)-value \\
\midrule
\(S_{\rm flux}\) & \(0.741\) & \(1.84\times10^{-4}\) \\
\(S_{\rm poisson}\) & \(0.741\) & \(1.84\times10^{-4}\) \\
\(\Phi_{\rm BuP}\) & \(0.738\) & \(2.01\times10^{-4}\) \\
\(|\Phi_{\rm BuP}|\) & \(-0.069\) & \(0.772\) \\
\(|\nabla\Phi_{\rm BuP}|\) & \(-0.582\) & \(0.00710\) \\
\bottomrule
\end{tabular}
\caption{
Couplage de la source émergente à une équation de Poisson discrète. Le potentiel
\(\Phi_{\rm BuP}\) reste fortement corrélé à la réponse de courbure.
}
\end{table}

\section{Scan en largeur de l'excitation}

Nous scannons ensuite :
\[
\sigma
=
0.01,\;0.02,\;0.05,\;0.08,\;0.10,\;0.15.
\]

\begin{table}[htbp]
\centering
\begin{tabular}{c c c l}
\toprule
\(\sigma\) & \(S_{\rm flux}\) vs \(|\delta R|\) & \(\Phi_{\rm BuP}\) vs \(|\delta R|\) & Lecture \\
\midrule
\(0.01\) & \(0.504\) & \(0.498\) & signal modéré \\
\(0.02\) & \(0.473\) & \(0.479\) & signal modéré \\
\(0.05\) & \(0.148\) & \(0.115\) & non significatif \\
\(0.08\) & \(0.176\) & \(0.157\) & non significatif \\
\(0.10\) & \(0.122\) & \(0.199\) & non significatif \\
\(0.15\) & \(0.741\) & \(0.738\) & meilleur régime \\
\bottomrule
\end{tabular}
\caption{
Scan en largeur d'excitation. Le meilleur couplage gravitationnel est obtenu
pour \(\sigma=0.15\).
}
\end{table}

Le résultat indique que le couplage gravitationnel optimal n'est pas obtenu pour
l'excitation la plus compacte, mais pour une excitation plus étendue. Cela
suggère que la gravité émergente BuP est un phénomène collectif du réseau
d'intrication.

\section{Structure du gradient}

Le gradient du potentiel vérifie au meilleur point :
\[
\rho_{\rm Spearman}(|\nabla\Phi_{\rm BuP}|,|\delta R|)
=
-0.582,
\qquad
p=0.00710.
\]

Cette anti-corrélation suggère que \(|\nabla\Phi_{\rm BuP}|\) ne trace pas le
pic central de courbure, mais plutôt une structure périphérique ou de bord
autour de l'excitation.

Le potentiel lui-même est donc le meilleur prédicteur direct de la réponse de
courbure, tandis que le gradient semble porter une information spatiale
complémentaire.

\section{Prédictions observables}

Paper~9 conduit à plusieurs prédictions.

\subsection{Potentiel effectif}

Une excitation locale d'intrication doit générer un potentiel effectif
\[
\Phi_{\rm BuP}
\]
corrélé à la réponse de courbure.

\subsection{Gravité tensorielle}

La gravité émergente BuP ne dépend pas seulement de \(T_{00}\), mais de la
structure complète :
\[
T_{00},
\qquad
T_{aa},
\qquad
T_{\rm grad},
\qquad
T_{0a}.
\]

Deux excitations ayant la même énergie informationnelle totale peuvent donc
produire des réponses gravitationnelles différentes si leurs flux ou contraintes
internes diffèrent.

\subsection{Largeur optimale}

Le couplage gravitationnel est maximal pour une excitation étendue :
\[
\sigma=0.15.
\]
Cela suggère que la gravité BuP est une réponse collective du réseau.

\subsection{Applications futures}

Le potentiel \(\Phi_{\rm BuP}\) pourra être utilisé pour construire :

\[
a_{\rm BuP}(r)=|\nabla\Phi_{\rm BuP}(r)|,
\]

des profils de rotation :

\[
v^2(r)\sim r\,\partial_r\Phi_{\rm BuP},
\]

ou des profils de lentille :

\[
\alpha_{\rm BuP}(r)\propto \nabla_\perp\Phi_{\rm BuP}.
\]

\section{Discussion}

Le résultat central de Paper~9 est que la source tensorielle de Paper~8 ne se
contente pas d'être corrélée directement à la courbure. Lorsqu'elle est
propagée par une équation de Poisson discrète sur le graphe d'intrication, elle
génère un potentiel effectif qui reste fortement corrélé à la réponse
géométrique.

Cela suggère l'existence d'une dynamique gravitationnelle effective interne au
réseau BuP :
\[
\text{source tensorielle}
\longrightarrow
\text{potentiel}
\longrightarrow
\text{courbure}.
\]

\section{Limites}

Cette étude reste un test numérique minimal. Le potentiel est construit par une
équation de Poisson discrète, et non encore par une résolution complète de
l'équation d'Einstein effective.

La normalisation physique de \(G_{\rm eff}(d_s)\), le rôle complet de
\(T_{\mu\nu}^{\rm ent}\), les effets de taille \(N\), et la comparaison avec des
observables astrophysiques réelles sont laissés à des travaux futurs.

\section{Conclusion}

Paper~9 couple la source matière émergente de Paper~8 à une équation de Poisson
discrète sur le graphe d'intrication :
\[
L_{\rm ent}\Phi_{\rm BuP}
=
S_{\rm flux}.
\]

Le résultat principal est :
\[
\rho_{\rm Spearman}(\Phi_{\rm BuP},|\delta R|)
=
0.738,
\qquad
p=2.01\times10^{-4}.
\]

Ainsi, le potentiel effectif BuP prédit fortement la réponse de courbure. Ce
résultat constitue une première prédiction gravitationnelle issue de la matière
émergente BuP.

\begin{thebibliography}{9}

\bibitem{hamdad-paper5}
F.~Hamdad,
\textit{Équation d'Einstein effective émergente dans la Gravité Quantique Bottom-Up},
BuP Paper 5, 2026.

\bibitem{hamdad-paper7}
F.~Hamdad,
\textit{Étude de la dynamique des excitations du réseau d'intrication BuP},
BuP Paper 7, 2026.

\bibitem{hamdad-paper8}
F.~Hamdad,
\textit{De l'excitation d'intrication à la matière émergente},
BuP Paper 8, 2026.

\bibitem{ollivier}
Y.~Ollivier,
\textit{Ricci curvature of Markov chains on metric spaces},
Journal of Functional Analysis, 256, 810--864, 2009.

\end{thebibliography}

\end{document}
